\chapter{Dimension Theory}

\begin{remark}
    We shall continue to assume that all rings,
    unless otherwise noted, are commutative.
\end{remark}

\section{Artinian Rings}

\begin{remark}
    Despite its superficial connection with Noetherian rings,
    the theory of Artinian rings is much simpler.
\end{remark}

\begin{proposition}\label{prop:artinian-prime-maximal}
    Prime ideals are maximal in Artinian rings.
\end{proposition}
\begin{proof}
    
\end{proof}
\begin{corollary}
    The nilradical is the same as the Jacobson radical in Artinian rings.
\end{corollary}
\begin{proof}
    The nilradical is the intersection of all prime ideals,
    and the Jacobson radical is the intersection of all maximal ideals,
    which we note by the \hyperref[prop:artinian-prime-maximal]{proposition above}
    to be the same object.
\end{proof}

\begin{proposition}
    Artinian rings have finitely many maximal ideals.
\end{proposition}
\begin{proof}
    
\end{proof}
\begin{proposition}
    The nilradical in Artinian rings is nilpotent.
\end{proposition}
\begin{proof}
    
\end{proof}

\begin{definition}
    Suppose \(R\) is a ring, and \(\mfk{p} \subseteq R\) a prime ideal.
    The height of \(\mfk{p}\) is
    \(\height\mfk{p} = \sup\{n: \mfk{p}_0 \subsetneq \mfk{p}_1 \subsetneq \cdots \subsetneq \mfk{p}_n = \mfk{p}\}\),
    the length of the longest strictly ascending chain of prime ideals terminating at \(\mfk{p}\).
    The Krull dimension of \(R\) is \(\dim R = \sup_{\mfk{p}\,\text{prime}}\height\mfk{p}\),
    the length of the longest strictly ascending chain of prime ideals.
\end{definition}

\begin{theorem}
    Suppose \(R\) is a ring.
    \(R\) is Artinian if and only if \(R\) is Noetherian with \(\dim R = 0\).
\end{theorem}
\begin{proof}
    
\end{proof}

\begin{proposition}
    Suppose \((R,\mfk{m})\) is a Noetherian local ring.
    Then \(\mfk{m}^n = 0\) for some \(n\) if and only if \(\mfk{m}^k = \mfk{m}^{k+1}\) for some \(k\).
    In that case, \(R\) is an Artinian local ring.
\end{proposition}
\begin{proof}
    
\end{proof}

\begin{theorem}[Structure Theorem for Artinian Rings]
    Any Artinian ring \(R\) can be uniquely written as
    a finite direct sum of Artinian local rings \(R = \bigoplus_{i=1}^r R_i\).
\end{theorem}
\begin{proof}
    
\end{proof}

\begin{proposition}
    Suppose \((R,\mfk{m})\) is an Artinian local ring.
    Then the following are equivalent:
    \begin{enumerate}[label={(\roman*)}, itemsep=0mm]
        \item \(R\) is a PID;
        \item \(\mfk{m}\) is a principal ideal; and
        \item the vector space dimension \(\dim_{R/\mfk{m}}(\mfk{m}/\mfk{m}^2) \leq 1\).
    \end{enumerate}
\end{proposition}
\begin{proof}
    
\end{proof}

\begin{proposition}
    Suppose \(K\) is a field, and \(R\) is finite-type over \(K\).
    \(R\) is Artinian if and only if \(R\) is finite over \(K\).
\end{proposition}
\begin{proof}
    
\end{proof}


\section{Discrete Valuation Rings}


\section{Gr\"obner Basis}
